\documentclass[UTF8,fontset=macnew,11pt]{ctexart}

\usepackage[a4paper,margin=2.2cm]{geometry}
\usepackage{amsmath,amssymb}
\usepackage{booktabs}
\usepackage{tabularx}
\usepackage{array}
\usepackage{float}
\usepackage{graphicx}
\usepackage{xcolor}
\usepackage{hyperref}

\hypersetup{
  colorlinks=true,
  linkcolor=blue!50!black,
  citecolor=blue!50!black,
  urlcolor=blue!50!black
}

\graphicspath{{../paper_latex/figures/submission_v36/}}
\setlength{\parskip}{0.34em}
\setlength{\parindent}{2em}
\setlength{\emergencystretch}{2em}

\newcommand{\method}{Resource-NMCTS}
\newcommand{\anf}{\ensuremath{\mathrm{ANF}}}
\newcommand{\fprm}{\ensuremath{\mathrm{FPRM}}}
\newcommand{\mcts}{\ensuremath{\mathrm{MCTS}}}
\newcommand{\xor}{\ensuremath{\oplus}}
\newcommand{\mct}{\ensuremath{\mathrm{MCT}}}
\newcommand{\rz}{\ensuremath{R_z}}
\newcommand{\score}{\ensuremath{\mathrm{score}}}
\newcommand{\ttable}{truth table}

\title{面向资源约束量子布尔函数综合的神经蒙特卡洛树搜索方法}
\author{中文论文稿 v36：图表化投稿主稿}
\date{2026年7月8日}

\begin{document}
\pagestyle{plain}
\maketitle

\begin{abstract}
量子布尔函数 oracle 综合是 Grover 搜索、约束求解和量子密码分析等算法中的基础编译环节。现有方法通常从固定代数正规形、ESOP/LUT/XAG 逻辑网络、可逆模板或 CNOT-oriented 空间结构出发；这些路线在单一指标或小规模函数上有效，但难以在 T-count、CNOT、深度、门数和辅助比特之间进行统一的资源约束搜索。本文把 bit-flip Boolean oracle synthesis 建模为 \anf{}/\fprm{} 项集合上的组合优化问题，并提出 \method{}，一种结合神经动作先验、蒙特卡洛树搜索、布尔环线性因子筛选、结构深度策略、frontier policy 与保守 guard 的逻辑层综合框架。本文只讨论逻辑层线路资源，不包含物理设备 mapping、连通性约束、噪声模型或 magic-state factory 调度。

在 $n\leq6$ 的 177 个传统函数上，Pareto-\method{} 相对 direct \anf{} 平均 T-count 降低 72.25\%、平均 score 降低 67.80\%；相对 ESOP beam、weighted ESOP-MILP 和 SSHR-H 分别获得 174/0/3、167/3/7 和 173/0/4 的配对胜/负/平。相对 ROS-style LUT proxy，Pareto-\method{} 在 309 个函数上为 309/0/0，平均 score 降低 84.27\%；相对官方 mockturtle KLUT-to-XAG probe，在传统 177 个函数上为 166/11/0，在 $n=14$ 的 64 个高维函数上为 64/0/0。内部消融显示，Affine-NMCTS 相对 fixed-coordinate MCTS 平均 score 降低 14.82\%，Pareto archive 相对单一 \method{} 平均 score 降低 3.26\%，learned prior 提供小幅但无 score-loss 的质量增益。

为处理 RevKit \texttt{oracle\_synth} 的 phase/Rz gate-set 边界，本文实现了可验证的 phase-parity \anf{} emitter、fixed-polarity \fprm{} search 和 Affine-FPRM phase search。Affine-FPRM 在 $T/R_z=30$ proxy 下相对 fixed-polarity FPRM 为 81/0/96、平均 score 降低 2.51\%，相对 RevKit 为 177/0/0、平均降低 65.50\%。正确性方面，本稿主图汇总的 $n=20$--$40$ item-level runs 为 5640/5640，width probe 为 1512/1512，phase selected rows 为 531/531；完整 \ttable{} bridge 在 $n=21$--$25$ 上累计 400/400 方法行通过完整 oracle、plan \anf{} 和 emitted-circuit 符号验证。结果表明，资源约束量子布尔函数综合可以被写成可验证的搜索问题；当前边界是 phase/Rz 分支仍为 parity-gadget 级 proxy，外部 ROS/legacy RevKit-CirKit 仍未完整复现，硬件 mapping 不在本文范围内。
\end{abstract}

\noindent\textbf{关键词：}量子 oracle 综合；布尔函数综合；强化学习；蒙特卡洛树搜索；ANF；FPRM；资源约束编译

\section{引言}

给定布尔函数 $f:\{0,1\}^n\rightarrow\{0,1\}$，Boolean oracle synthesis 需要生成实现
\begin{equation}
  |x\rangle |y\rangle \mapsto |x\rangle |y\xor f(x)\rangle
  \label{eq:oracle}
\end{equation}
的可逆逻辑线路。这个任务看似只是把经典逻辑嵌入量子程序，但在 fault-tolerant 或资源受限量子计算场景中，布尔表示、计算/反计算顺序、辅助比特复用和相位实现方式都会显著改变 T-count、CNOT、深度和辅助比特需求。因此，oracle synthesis 不是一个单步翻译问题，而是一个带有多资源目标和语义约束的离散搜索问题。

已有工作从多个方向处理这个问题。ESOP、XAG、LUT 和 BDD 类方法通过经典逻辑网络降低乘法复杂度或局部函数成本\cite{wille2009bdd,meuli2020ros,meuli2022xag}；RevKit 和 mockturtle 提供了量子/经典逻辑综合工具链入口\cite{soeken2018epfl,revkit,mockturtle}；SSHR 则利用 hypercube 中 parallelotope 空间结构进行 small-scale CNOT-oriented synthesis\cite{zheng2025sshr}。这些方法给出了重要基线，但它们通常把候选结构或目标指标固定在某一类表示中。对于同时关注 T-count、CNOT、depth、gate count 和 peak ancilla 的逻辑层 oracle 综合，固定表示不一定能覆盖足够的资源折中空间。

本文的出发点是把 oracle synthesis 的主要自由度显式写成搜索空间。我们用 \anf{} 与 \fprm{} 项集合表示状态，用布尔环因子、极性选择、线性奇偶因子和可逆仿射变换定义动作，用 \mcts{} 与神经先验分配搜索预算，并用 guard 保证学习候选不破坏确定性 baseline。这个设计使 AI 不直接生成无约束门序列，而是在可验证的代数动作之间排序。最终线路可以通过 \anf{} 展开、GF(2) 符号模拟和完整 \ttable{} 逐点验证进行检查。

本文贡献如下：
\begin{itemize}
  \item 提出 \method{}，把量子布尔函数 oracle 综合建模为 \anf{}/\fprm{} 项集合上的资源约束搜索问题，而不是 SSHR parallelotope 空间的局部扩展。
  \item 设计布尔环线性因子筛选、神经动作先验、\mcts{}、depth-frontier policy、Pareto archive 和 guard 组成的逻辑层 synthesis pipeline。
  \item 在 ESOP、SSHR、ROS-style LUT proxy、mockturtle XAG probe 和 RevKit phase/Rz proxy 上给出配对对比，并用内部消融拆分搜索、学习和 portfolio 的贡献。
  \item 用 plan \anf{} 验证、emitted-circuit GF(2) 符号验证和 $n=21$--$25$ 完整 \ttable{} bridge 构建正确性证据链。
\end{itemize}

\begin{figure}[H]
\centering
\includegraphics[width=\linewidth]{fig1_pipeline.pdf}
\caption{\textbf{\method{} 的可验证搜索 pipeline。} 输入可以是完整 \ttable{} 或 \anf{} 项集合；动作空间包含布尔环因子、极性、线性奇偶因子和 affine phase 分支；搜索由神经先验、\mcts{} 和 frontier policy 控制；最终输出在 plan、emitted circuit 和完整 \ttable{} 层进行验证。Source data: \texttt{fig1\_pipeline.csv}.}
\label{fig:pipeline}
\end{figure}

\section{相关工作}

\subsection{布尔表示与 oracle 综合}

低 T-count oracle 综合长期依赖合适的布尔中间表示。BDD-based reversible synthesis 以符号决策图处理大函数\cite{wille2009bdd}；ROS 把 oracle synthesis 写成 resource-constrained LUT mapping，并强调 LUT 分解和后续可逆实现的资源耦合\cite{meuli2020ros}；XAG 将 XOR 与 AND 结构分离，使 AND 节点数成为非 Clifford 成本的重要 proxy\cite{meuli2022xag}。mockturtle 进一步提供了可复用的 KLUT、XAG 和重综合算法库\cite{mockturtle}。这些路线说明，经典逻辑综合表示会直接影响量子 oracle 资源。

本文与上述工作共享“布尔表示决定量子资源”的前提，但不把搜索限制在某一种固定网络中。我们的状态是 \anf{}/\fprm{} 项集合，动作是可展开回 GF(2) 多项式的代数分解。这样做牺牲了部分已有工具链的成熟优化，但换来了更直接的语义验证和更自然的多资源候选组合。

\subsection{SSHR 与小规模空间结构方法}

SSHR 利用 Boolean hypercube 中 parallelotope 结构生成 MCT block，并以 CNOT reduction 为主要目标\cite{zheng2025sshr}。它适合本文作为 small-scale CNOT-oriented baseline，因为它关注小规模布尔函数、CNOT 计数和空间结构候选。本文不把 SSHR 当作主线方法：\method{} 的状态、动作和验证均定义在 \anf{}/\fprm{} 项集合和代数搜索上。因此，本文可以在结果中承认 SSHR-H 的 CNOT 优势，同时主张自身在 T-count 和加权 score 上更强。

\subsection{学习式搜索与线路综合}

学习式搜索已经被用于经典电路设计、量子线路综合和图重写优化。ShortCircuit 使用 AlphaZero 风格搜索处理经典布尔电路设计\cite{tsaras2024shortcircuit}；Gumbel AlphaZero 被用于 Clifford+T unitary synthesis\cite{rietsch2024unitary}；强化学习也被用于 ZX-calculus rewrite-rule selection\cite{riu2025rlzx}；MonteQ 使用 \mcts{} 处理量子线路合成中的动作排序\cite{machiya2026monteq}。这些工作共同表明，离散综合问题的动作空间通常过大，需要学习策略或树搜索分配有限预算。

本文的区别在于，学习策略不直接写门序列，而是对可验证的代数动作排序。正确性来自代数展开和 circuit-level symbolic simulation，而不是来自模型预测本身。这使得 learned prior 的贡献可以通过消融量化；在当前数据中，它是小幅正向模块，而非最大增益来源。

\section{问题定义与资源模型}

本文考虑 bit-flip Boolean oracle。输入函数可以表示为完整 \ttable{}，也可以表示为 square-free \anf{}
\begin{equation}
  f(x)=\bigoplus_{m\in\mathcal{M}}\prod_{i\in m}x_i ,
  \label{eq:anf}
\end{equation}
其中 $\mathcal{M}$ 是 monomial 集合。输出线路使用 X、CNOT 和多控 Toffoli（\mct{}）抽象门，并在逻辑层估计 T-count、CNOT、depth、gate count 与 peak ancilla。

搜索排序采用统一加权分数
\begin{equation}
  \score = T + 0.04\,\mathrm{CNOT} + 0.015\,\mathrm{depth}
         + 0.01\,\mathrm{gates} + 2\,\mathrm{ancilla}.
  \label{eq:score}
\end{equation}
这个分数只用于逻辑层候选排序，并不代表真实设备运行时间或物理错误率。因此，本文同时报告 T、CNOT、depth 和 ancilla，避免单一 score 掩盖 trade-off。

验证分为三层。第一层在小规模函数和高维 bridge 中构造完整 \ttable{}，逐点验证式~\eqref{eq:oracle}。第二层把 factor plan 展开回 \anf{} 项集合，检查目标多项式是否一致。第三层对 emitted X/CNOT/\mct{} circuit 做 GF(2) 符号模拟，检查输出线多项式、输入线复原和辅助线清零。第三层不枚举 $2^n$ 个输入，但验证的是最终 emitted circuit，而不只是搜索 plan。

\section{方法}

\subsection{项集合状态与可验证动作}

\method{} 的状态是当前尚未综合的项集合 $\mathcal{M}$。一次动作选择一个可计算、可反算的代数结构，将 $\mathcal{M}$ 分解为 quotient 子问题和 rest 子问题，再递归生成 compute/action/uncompute 线路。候选动作包括公共单项式因子、\fprm{} 极性重写后的公共项、线性奇偶因子、布尔环线性因子和 root-action 扩展。动作被接受后，框架估计即时资源收益、子问题大小、辅助比特峰值和后续可分解性。

神经动作先验用于对候选排序，\mcts{} 在有限预算内探索候选组合，rollout 使用确定性资源估计和已有 baseline 候选。最终 portfolio 在多个候选线路中按式~\eqref{eq:score} 选择最优者。这里 AI 的职责是搜索控制，而不是语义生成；任何候选线路都必须通过 \anf{} 或 circuit-level 验证。

\subsection{布尔环线性因子}

传统线性因子筛选往往要求线性因子变量和 quotient 变量不相交。本文的 Boolean-ring screen 允许二者共享变量，并在展开时使用 $x_i^2=x_i$ 与 GF(2) 偶数项消去。例如
\begin{equation}
  (x_i\xor x_j)x_i m = x_i m \xor x_i x_j m .
  \label{eq:boolean-ring}
\end{equation}
线路层面仍可用 CNOT 计算线性奇偶量，再以该辅助量控制 quotient 子线路。该设计扩大了可搜索代数结构，而不是单纯增加 beam width。

\subsection{Frontier policy、Pareto archive 与 guard}

本文把 AI 模块拆为三类。第一类是动作排序，模型根据候选覆盖率、项数变化、变量覆盖密度、次数分布和即时资源估计对候选动作打分。第二类是结构选择，depth policy 判断当前项集合适合 single、depth-1 还是 depth-2 screen。第三类是 frontier 选择，depth-frontier policy 在 depth-2、depth-3 和 depth-4 screen 之间选择；quality policy 接近 all-depth frontier，cost-aware policy 用小幅 score 代价换取更低 plan time 和辅助线 lifetime area。

Guard 负责限制学习策略风险。若学习候选相对确定性 baseline 出现 score 退化，则返回 baseline。Pareto archive 在多个候选中保留不同资源折中，再由目标 score 选择最终输出。这个设计使算法可以报告“质量提升”“时间节省”和“质量-时间折中”三类贡献。

\section{实验设置}

实验回答四个问题。第一，\method{} 在完整 \ttable{} 可验证的小规模函数上是否降低 T-count 和 score。第二，与 ESOP、SSHR、ROS-style LUT proxy、mockturtle XAG probe 和 RevKit phase/Rz proxy 相比，结论是否仍成立。第三，内部消融是否能证明提升来自 affine search、\mcts{}、learned prior、guard 和 Pareto archive，而不是只来自换计分口径。第四，高维项集和 bridge 是否具有足够强的语义验证证据。

传统函数切片包含 $n\leq6$ 的 177 个函数，baseline 包括 direct \anf{}、AND-direct \anf{}、ESOP beam、weighted ESOP-MILP、SSHR-H、ROS-style LUT proxy、mockturtle KLUT-to-XAG probe 和 RevKit Python \texttt{oracle\_synth}。高维切片包含 $n=14$ mockturtle probe、$n=18,20$ 函数级边界、$n=20$--$40$ 项集级 scale、width probe，以及 $n=21,22,23,24,25$ 完整 \ttable{} bridge。配对比较只纳入函数名或项集名匹配、语义正确且未 timeout 的结果。

\section{结果}

\subsection{传统函数上的多资源优势}

图~\ref{fig:traditional} 和表~\ref{tab:traditional} 给出 $n\leq6$ 传统函数切片上的平均资源。Pareto-\method{} 相对 direct \anf{} 平均 T-count 降低 72.25\%、平均 score 降低 67.80\%；相对 weighted ESOP-MILP 平均 T-count 降低 32.77\%、score 降低 29.84\%。SSHR-H 的平均 CNOT 最低，说明本文不能声称 CNOT-only 支配 SSHR；更准确的结论是，本文以一定 CNOT 和辅助比特 trade-off 换取更低 T-count、depth 和加权 score。

\begin{figure}[H]
\centering
\includegraphics[width=\linewidth]{fig2_traditional_resources.pdf}
\caption{\textbf{传统函数上的平均资源。} Pareto-\method{} 在 T-count 和 weighted score 上最低；SSHR-H 仍是强 CNOT-oriented baseline。Source data: \texttt{fig2\_traditional\_resources.csv}.}
\label{fig:traditional}
\end{figure}

\begin{table}[H]
\centering
\small
\caption{$n\leq6$ 传统函数切片上的平均逻辑资源。}
\label{tab:traditional}
\begin{tabular}{lrrrrr}
\toprule
方法 & 平均 T & 平均 CNOT & 平均 depth & 平均 ancilla & 平均 score \\
\midrule
Direct ANF & 184.1 & 134.2 & 134.6 & 1.33 & 194.3 \\
AND-direct ANF & 101.0 & 166.5 & 166.9 & 2.18 & 115.2 \\
ESOP-MILP & 83.6 & 133.5 & 159.6 & 2.32 & 96.7 \\
SSHR-H & 81.0 & \textbf{67.6} & 88.7 & 1.36 & 88.2 \\
\method{} & 43.9 & 89.9 & 94.5 & 1.92 & 53.2 \\
Pareto-\method{} & \textbf{40.4} & 83.0 & \textbf{88.0} & 2.03 & \textbf{49.6} \\
\bottomrule
\end{tabular}
\end{table}

\subsection{外部 baseline 与可比边界}

图~\ref{fig:baselines} 汇总外部和传统 baseline 的配对比较。Pareto-\method{} 相对 ESOP beam 为 174/0/3，平均 score 降低 36.09\%；相对 ESOP-MILP 为 167/3/7，平均 score 降低 29.84\%；相对 SSHR-H 的 T-count 为 173/0/4，平均降低 47.93\%。ROS-style LUT proxy 覆盖 309 个函数，best-$K$ proxy 比固定 $K=4$ 更强，但 Pareto-\method{} 相对该 proxy 仍为 309/0/0。mockturtle probe 用官方 header 将 KLUT network 重综合为 XAG；在传统 177 个函数上为 166/11/0，在 $n=14$ 的 64 个函数上为 64/0/0。

\begin{figure}[H]
\centering
\includegraphics[width=0.94\linewidth]{fig3_baseline_comparisons.pdf}
\caption{\textbf{配对基线比较。} 横轴为目标方法相对基线的平均变化；负值表示资源更低。数字标签给出胜/负/平；源数据见 \texttt{fig3\_baseline\_comparisons.csv}.}
\label{fig:baselines}
\end{figure}

这些对比有意义，但需要严格限定。ROS-style LUT 是 proxy，不是官方 ROS 全流程；mockturtle probe 是 KLUT-to-XAG 逻辑网络比较，不输出可逆 Clifford+T 线路；RevKit \texttt{oracle\_synth} 返回 Rz-phase netlist，不能直接当作精确 T/Tdg 网络比较。本文据此只主张逻辑层资源优势，而不是工具链全流程或硬件后端优势。

\subsection{内部消融：提升来自哪里}

表~\ref{tab:ablation} 显示，最大算法增益来自 affine search。Full Affine-NMCTS 相对 fixed-coordinate MCTS 的平均 score 降低 14.82\%；Resource-NMCTS 相对 Affine-NMCTS 降低 1.91\%；Pareto archive 相对单一 \method{} 降低 3.26\%。Learned prior 对 \method{} 的 score 增益为 1.10\%，且无 score-loss，但运行时间上并非最优。因此，本文不把“神经网络”写成决定性单点突破，而把它定位为可验证搜索框架中的小幅正向搜索先验。

\begin{table}[H]
\centering
\scriptsize
\caption{搜索贡献分解。负的 mean $\Delta$ score 表示目标方法更优。}
\label{tab:ablation}
\begin{tabular}{@{}p{0.40\linewidth}p{0.20\linewidth}rcc@{}}
\toprule
Component & Dataset & Pairs & Score W/L/T & Mean $\Delta$ score \\
\midrule
Affine greedy vs fixed-coordinate MCTS & ablation\_affine & 321 & 165/88/68 & $-12.12\%$ \\
Neural refine over affine greedy & ablation\_affine & 322 & 65/0/257 & $-1.08\%$ \\
Guarded Affine-NMCTS over no-guard & ablation\_affine & 322 & 88/0/234 & $-1.74\%$ \\
Learned prior for Resource-NMCTS & traditional\_resource & 177 & 39/0/138 & $-1.10\%$ \\
Pareto archive over Resource-NMCTS & traditional\_resource & 177 & 68/0/109 & $-3.26\%$ \\
No-MCTS portfolio over heuristic-only & trad. ablation & 177 & 54/0/123 & $-2.51\%$ \\
Resource-NMCTS over no-MCTS portfolio & trad. ablation & 177 & 54/0/123 & $-1.44\%$ \\
Pareto Resource-NMCTS over no-MCTS portfolio & trad. ablation & 177 & 106/0/71 & $-4.69\%$ \\
Highdim no-MCTS portfolio over root beam & n14 ablation & 16 & 14/0/2 & $-6.25\%$ \\
Highdim no-MCTS portfolio over linear-pair & n14 ablation & 16 & 14/0/2 & $-3.08\%$ \\
Linear-pair guard vs root beam at n=14 & highdim\_resource & 64 & 60/0/4 & $-3.00\%$ \\
Recursive linear-pair guard vs root beam at n=15 & n15 scale & 32 & 30/0/2 & $-5.28\%$ \\
Recursive linear-pair guard vs shallow linear-pair at n=16 & n16 scale & 24 & 23/0/1 & $-2.54\%$ \\
Recursive linear-pair guard vs root beam at n=16 & n16 scale & 24 & 23/0/1 & $-4.31\%$ \\
Fast linear-pair guard vs root beam at n=18 & n18 stress & 12 & 6/0/6 & $-1.91\%$ \\
Resource-NMCTS vs fast linear-pair at n=18 & n18 stress & 12 & 12/0/0 & $-3.55\%$ \\
\bottomrule
\end{tabular}

\end{table}

\subsection{Phase/Rz 边界与 Affine-FPRM 搜索}

RevKit \texttt{oracle\_synth} 在 177 个传统函数上全部返回，但 171/177 行含非 Clifford \rz{}，总数为 9242。因此，如果只按 lower-bound phase netlist 计分，比较会低估 RevKit 输出在 Clifford+T 近似综合中的成本。为把该边界转化为本文内部可验证分支，我们实现 phase-parity \anf{} emitter，并进一步加入 fixed-polarity \fprm{} search 和 Affine-FPRM phase search。

图~\ref{fig:phase} 显示，Affine-FPRM 在 $T/R_z=30$ proxy 下相对 fixed-polarity FPRM 为 81/0/96，平均 score 降低 2.51\%；相对 phase-parity \anf{} 为 85/0/92，平均降低 2.98\%；相对 RevKit 为 177/0/0，平均降低 65.50\%。在 177 个函数中，81 个选择非 identity 线性变换，43 个选择非零极性。这个结果说明 phase/Rz 分支已经从成本敏感性讨论推进为可验证的代数搜索；但它仍是 parity-gadget 级逻辑 emitter，不包含近似旋转序列生成或硬件映射。

\begin{figure}[H]
\centering
\includegraphics[width=\linewidth]{fig4_phase_affine.pdf}
\caption{\textbf{Affine-FPRM 相位搜索。} 左图为 Affine-FPRM 相对 fixed FPRM、ANF 和 RevKit proxy 的平均变化；右图为非 identity 线性变换和非零极性的选择频率；源数据见 \texttt{fig4\_phase\_affine.csv}.}
\label{fig:phase}
\end{figure}

\begin{table}[H]
\centering
\small
\caption{Affine-FPRM phase-parity search 的配对比较。}
\label{tab:phase}
\begin{tabular}{lrrr}
\toprule
Comparison & Items & W/L/T & Mean $\Delta$ \\
\midrule
Affine-FPRM-$T/R_z=30$ vs FPRM & 177 & 81/0/96 & -2.51\% \\
Affine-FPRM-$T/R_z=30$ vs ANF & 177 & 85/0/92 & -2.98\% \\
Affine-FPRM-$T/R_z=30$ vs RevKit & 177 & 177/0/0 & -65.50\% \\
Affine-FPRM non-Clifford Rz vs FPRM & 177 & 20/0/157 & -0.73\% \\
Affine-FPRM-score vs FPRM & 177 & 80/0/97 & -5.77\% \\
\bottomrule
\end{tabular}

\end{table}

\subsection{高维验证与 bridge}

图~\ref{fig:validation} 汇总本稿主图采用的验证覆盖。$n=20$--$40$ item-level runs 为 5640/5640，width probe 为 1512/1512，phase selected rows 为 531/531，完整 \ttable{} bridge 为 400/400。Bridge 的组成是：$n=21,22$ 基础 bridge 120 行，$n=23$ 基础 bridge 40 行，$n=23$ large-policy rerun 60 行，$n=23$ cost-aware rerun 60 行，$n=24$ 和 $n=25$ 各 60 行。所有这些方法行同时通过完整 oracle、plan \anf{} 和 emitted-circuit 符号验证。

\begin{figure}[H]
\centering
\includegraphics[width=0.94\linewidth]{fig5_validation.pdf}
\caption{\textbf{验证覆盖。} 主图汇总的逻辑层结论均有对应 source-data 行；完整 \ttable{} bridge 当前按源数据为 400/400，而不是旧稿中曾使用的 420/420 口径。Source data: \texttt{fig5\_validation.csv}.}
\label{fig:validation}
\end{figure}

\begin{table}[H]
\centering
\small
\caption{高维验证与边界。}
\label{tab:validation}
\begin{tabularx}{\linewidth}{>{\raggedright\arraybackslash}p{0.34\linewidth}>{\raggedright\arraybackslash}p{0.22\linewidth}>{\raggedright\arraybackslash}X}
\toprule
验证项 & 当前覆盖 & 边界 \\
\midrule
$n=20$--$40$ item-level symbolic runs & 5640/5640 & 不是完整 \ttable{} 枚举。 \\
Width probe symbolic runs & 1512/1512 & 证明加宽候选不是主收益来源，但仍是项集级验证。 \\
$n=21$--$25$ complete bridge & 400/400 & 小样本完整枚举，不能外推为 $n=26$--$40$ 全部完整验证。 \\
Phase selected rows & 531/531 & up to global phase 验证，不输出近似旋转门序列。 \\
\bottomrule
\end{tabularx}
\end{table}

\section{讨论}

本文的核心提升不在于“比 SSHR 多一个启发式”，而在于把 oracle synthesis 改写成资源约束搜索问题。SSHR 在本文中是重要的 CNOT-oriented small-scale baseline；ESOP、ROS-style LUT、mockturtle XAG 和 RevKit phase netlist 则分别从乘积项、LUT、逻辑网络和 phase/Rz gate-set 角度提供外部压力测试。这样设置对比对象是有意义的，因为它们分别覆盖了本文可能被审稿人质疑的四个问题：是否只赢 naive ANF，是否只赢 SSHR，是否能对齐外部工具链，是否处理 RevKit phase/Rz 边界。

同时，当前稿件仍有清晰边界。第一，learned prior 的增益小于 affine search 和 Pareto archive，不能夸大为“深度强化学习主导性能提升”。第二，RevKit 和 ROS 的比较仍包含 proxy 或 adapter 成分，后续若要投稿到更严格的量子编译会议，应继续复现官方 ROS 或 legacy RevKit/CirKit 完整流程。第三，phase/Rz 分支虽然已有 Affine-FPRM 搜索，但仍停留在 parity-gadget 级 proxy；下一步应加入 learned phase/Rz policy、旋转谱优化和实际近似旋转序列审计。第四，所有结果都是逻辑层资源，不包含 routing、native gate mapping 或设备噪声。

因此，本文当前最稳妥的投稿定位是：一种逻辑层、资源约束、可验证的神经 \mcts{} Boolean oracle synthesis 方法。它可以声称在 T-count 和 weighted score 上显著优于多类 baseline，并在高维项集上有强符号验证和小样本完整 bridge；它不能声称硬件级最优、CNOT-only 支配 SSHR、或 phase/Rz 全局最优。

\section{结论}

本文提出 \method{}，将量子布尔函数 oracle 综合建模为 \anf{}/\fprm{} 项集合上的资源约束搜索问题，并结合神经动作先验、\mcts{}、布尔环线性因子筛选、frontier policy、Pareto archive 和 guard 生成逻辑层可逆线路。实验显示，Pareto-\method{} 在传统函数上相对 ANF、ESOP、SSHR-H 和 ESOP-MILP 获得稳定低 T-count 与低 score 优势，在 ROS-style LUT proxy 和 official-header mockturtle XAG probe 下仍保持优势。内部消融说明主要增益来自 affine search、资源感知 search 和 Pareto archive，learned prior 是小幅正向模块。Affine-FPRM phase search 进一步把 RevKit/Rz 边界推进到可验证的 phase oracle 搜索分支。

当前 v36 已经从实验日志推进为图表化投稿主稿，但目标尚未完全结束：更强的 learned phase/Rz-aware search、官方 ROS 或 legacy RevKit/CirKit 完整复现，以及面向目标期刊模板的英文投稿稿仍是下一阶段必须完成的工作。

\begin{thebibliography}{99}

\bibitem{wille2009bdd}
R.~Wille and R.~Drechsler.
\newblock BDD-based synthesis of reversible logic for large functions.
\newblock In \emph{Proceedings of DAC}, pp.~270--275, 2009.

\bibitem{meuli2020ros}
G.~Meuli, M.~Soeken, M.~Roetteler, and G.~De Micheli.
\newblock ROS: Resource-constrained oracle synthesis for quantum computers.
\newblock \emph{Electronic Proceedings in Theoretical Computer Science}, 318:119--130, 2020.

\bibitem{meuli2022xag}
G.~Meuli, M.~Soeken, and G.~De Micheli.
\newblock Xor-And-Inverter Graphs for quantum compilation.
\newblock \emph{npj Quantum Information}, 8:7, 2022.

\bibitem{soeken2018epfl}
M.~Soeken, G.~Meuli, E.~Testa, H.~Nada, M.~Reymond, B.~Dang, G.~De Micheli, and A.~Mishchenko.
\newblock The EPFL logic synthesis libraries.
\newblock \emph{arXiv:1805.05121}, 2018.

\bibitem{revkit}
M.~Soeken and collaborators.
\newblock RevKit: Python quantum compilation library.
\newblock \url{https://github.com/msoeken/revkit}.

\bibitem{mockturtle}
M.~Soeken and collaborators.
\newblock mockturtle: logic network optimization and synthesis.
\newblock \url{https://mockturtle.readthedocs.io/}.

\bibitem{zheng2025sshr}
Q.~Zheng, Y.~Xu, J.~Zhang, Z.~Su, and S.~Zheng.
\newblock CNOT oriented synthesis for small-scale Boolean functions using spatial structures of parallelotopes.
\newblock In \emph{Proceedings of ICCAD}, 2025.

\bibitem{tsaras2024shortcircuit}
D.~Tsaras, A.~Grosnit, L.~Chen, Z.~Xie, H.~Bou-Ammar, and M.~Yuan.
\newblock ShortCircuit: AlphaZero-driven circuit design.
\newblock \emph{arXiv:2408.09858}, 2024.

\bibitem{rietsch2024unitary}
S.~Rietsch, A.~Y. Dubey, C.~Ufrecht, M.~Periyasamy, A.~Plinge, C.~Mutschler, and D.~D. Scherer.
\newblock Unitary synthesis of Clifford+T circuits with reinforcement learning.
\newblock In \emph{IEEE International Conference on Quantum Computing and Engineering}, pp.~824--835, 2024.

\bibitem{riu2025rlzx}
J.~Riu, J.~Nogu{\'e}, G.~Vilaplana, A.~Garcia-Saez, and M.~P. Estarellas.
\newblock Reinforcement learning based quantum circuit optimization via ZX-calculus.
\newblock \emph{Quantum}, 9:1758, 2025.

\bibitem{machiya2026monteq}
M.~Machiya, M.~Menickelly, P.~Hovland, and J.~Liu.
\newblock MonteQ: A Monte Carlo tree search based quantum circuit synthesis framework.
\newblock \emph{arXiv:2604.19029}, 2026.

\end{thebibliography}

\end{document}
